\section{Introducción}

Riichi Mahjong es un juego de información imperfecta,
caracterizado por una elevada complejidad estratégica y por la presencia de
información oculta. Su espacio de estados se estima en el orden de $10^{48}$\cite{li2020suphx}.
Además, el orden de juego puede alterarse, ya que ciertas acciones, como el
\textit{pon}, el \textit{chi} o el \textit{kan}, permiten interrumpir la secuencia habitual de turnos. Esto incrementa aún más la complejidad del juego.

Este proyecto desarrolla una Inteligencia Artificial que predice una acción de descarte a partir de la información observable de un estado de juego,
que incluye información secuencial y estática. El modelo se basa en la arquitectura Transformer
y se entrena mediante aprendizaje por imitación utilizando 500.000 estados de juego reales
obtenidos de partidas de la plataforma Tenhou.net.

Durante el desarrollo del proyecto se obtienen diferentes métricas de precisión,
analizando tanto la coincidencia con el descarte observado como la precisión estratégica, que se basa en categorías semánticas propias del juego
calculadas mediante el \textit{Shanten} y el \textit{Ukeire}.

El esquema general de esta memoria comienza con
los objetivos del proyecto,
seguido de una sección de teoría, donde se explican
el juego de Mahjong y las heurísticas del juego,
así como la arquitectura de Transformers y el aprendizaje
por imitación.
Después, se realiza un repaso del estado del arte,
donde se analizan los trabajos relacionados con el juego de Mahjong y con el uso de Transformers en juegos de información imperfecta.
A continuación, se describen las herramientas utilizadas para el desarrollo del proyecto
y se explica su proceso de desarrollo, desde la preparación de los datos hasta el
entrenamiento del modelo y la evaluación de los resultados.
Finalmente, se presentan los resultados obtenidos y se discuten las conclusiones del proyecto.

El trabajo surgió inicialmente como un proyecto personal orientado a crear una IA de Mahjong que sirviese como apoyo para aprender a jugar mejor. A partir de esta idea inicial, el proyecto evolucionó hacia un trabajo académico de carácter exploratorio. Al comienzo no existía una arquitectura final ni un conjunto de herramientas previamente definido, por lo que las decisiones se tomaron de forma progresiva a partir de las pruebas realizadas y de los resultados obtenidos.

No se siguió una metodología formal concreta ni se estableció un calendario cerrado de hitos. El desarrollo fue iterativo e incremental: se probaron diferentes formas de representar los estados de juego, incluyendo una aproximación inicial mediante redes convolucionales, antes de adoptar una representación basada en tokens y una arquitectura Transformer. Esta solución resultó más natural para codificar los distintos elementos de un estado de Mahjong. De manera similar, la evaluación estratégica basada en Shanten y Ukeire no formaba parte del planteamiento inicial, sino que surgió durante el desarrollo al observar las limitaciones de la precisión exacta. Los cuadernos de Jupyter documentan estas iteraciones y permiten seguir la evolución del proyecto desde las primeras pruebas hasta los resultados finales.
